\paragraph{Recommendations.} For anyone who relabels a benchmark with model assistance, our results suggest the following.
\begin{enumerate}
\item \emph{Separate the goals.} Adjudicating flagged items cleans a benchmark; accuracy claims need a probability sample that includes concordant items.
\item \emph{Sample concordant items and report the sample}: the flagging rule, $|\Dset|$, the inclusion probabilities, and the concordant items adjudicated with the errors found. The difference estimator of Section~\ref{sec:twophase} is then valid for every model, now or later.
\item \emph{When label errors are rare, report the simultaneous intervals of Theorem~\ref{thm:simultaneous}}, which cover all accuracies, differences and ranking orders; Wald intervals are unreliable there (Section~\ref{sec:emp-sim}).
\item \emph{Do not rank the flagger, or models close to it, on its own audit without correction} (Corollary~\ref{cor:advantage}).
\item \emph{Re-analyse past DR audits with the sensitivity bounds of Theorem~\ref{thm:ident}}, asking how large a concordant error rate would overturn a reported ranking.
\item \emph{Allocate the budget by stratum error rates}: when estimation is the goal, adjudicating every flagged item first is inefficient (Table~\ref{tab:allocation}).
\end{enumerate}

\paragraph{Limitations.} Our reference labels are themselves the output of an annotation process. MMLU-Redux annotators can err, and items with ambiguous or multiple answers were excluded rather than modelled. All results are therefore relative to the adjudication process (Remark~\ref{rem:adjudicator}), as they would be in any real audit.

Scenario A has only \nLabelErrA{} label errors among \nItems{} items, so its estimates of hidden-error counts rest on few events. Scenario B is semi-synthetic: its labels come from a real model's answers to real questions, but no published benchmark was built this way from these models. The HELM predictions come from 5-shot prompting of models released in 2023--2024; stronger or differently prompted models would change the numbers but not the identities.

For other models the rate of the lower bounds depends on how often they repeat the hidden errors. Our simultaneous intervals are valid but certify few adjacent pairs when models are separated by fractions of a point; certifying such rankings needs large concordant samples. Both lower bounds are stated for the flagger's accuracy, with constants that are not tight. The Neyman allocation in Table~\ref{tab:allocation} uses oracle stratum error rates. Our evidence comes from simulated audits of one benchmark and from one published audit. In the published audit (CoNLL-2003) the hidden-error effect on accuracy is small, the reference re-annotation corrected existing labels rather than labelling blind, and the two expert relabelings disagree on many tokens; other benchmarks may show larger or smaller effects. Finally, our literature survey covers 22 studies chosen to span the design space, not a systematic review.

\paragraph{Conclusion.} Using models to decide which benchmark labels to check is efficient for finding errors, and it is here to stay. The accuracies computed afterwards, however, are not estimates of true accuracy. They inflate the flagger by exactly the mass of errors it shares with the benchmark, they can reorder rankings, and they lock in the blind spots of whichever models did the cleaning. A random sample of the items the models agreed on removes all three problems. When hidden errors are common it does so at a rate no design can beat. When they are rare, no design can do much better than DR's own error, but a concordant sample still reveals how large that error might be.
